\begin{enbref}
{\bfseries\color{bleu}L'ESSENTIEL}\par\vspace{3pt}
\begin{itemize}[leftmargin=12pt,itemsep=2pt,topsep=0pt]
\item \textbf{Le nouvel impôt sur le revenu élargit l'assiette de 113~Md€ et rend tout aux contribuables.} Il remplace
l'impôt sur le revenu, la CSG et la CRDS. Les niches n'apportent que 19~Md€ de cet élargissement, l'essentiel vient de
l'imposition des pensions comme des salaires, de l'impôt individuel et de la fin des déductions forfaitaires. Le taux
marginal maximal passe de 55~\% à 48~\%.
\item \textbf{L'impôt immobilier quitte les transactions pour le sol.} Les droits de mutation sont divisés par deux
($-$8,1~Md€) et une taxe nationale sur la valeur du sol rapporte 15~Md€.
\item \textbf{La déduction immédiate des investissements coûte cher pendant sa montée en charge.} Une fois retirés les
dispositifs déjà comptés dans la partie III, dont le recentrage du crédit d'impôt recherche, le bloc de l'impôt sur les
sociétés coûte 5,7~Md€ en 2032 et s'équilibre presque à terme.
\item \textbf{Le solde public est de $-$0,6~Md€ en 2032 et de +5,3~Md€ à terme.} Les prélèvements obligatoires
augmentent de 0,2 point de PIB à terme, marge qui peut financer d'autres baisses d'impôt.
\item \textbf{Trois arbitrages déplacent ces chiffres.} Supprimer aussi le tiers restant du crédit d'impôt recherche
rapporterait 2,1~Md€. Ramener le prélèvement forfaitaire unique à 25~\% coûterait 3,8~Md€. Chaque 0,05 point de taxe sur
le sol rapporte environ 3~Md€.
\end{itemize}
\end{enbref}

\section{1. Philosophie des réformes}

\textbf{Chaque niche fiscale oblige à taxer davantage ceux qui n'en bénéficient pas.} Notre système fiscal procède
souvent en deux temps. Il affiche d'abord des taux élevés, puis rend une partie de l'argent à ceux qui entrent dans la
bonne case : le ménage qui emploie une aide à domicile, l'épargnant qui choisit le bon placement, le retraité plutôt que
le salarié, le couple marié plutôt que le couple en concubinage, la petite entreprise plutôt que la grande. Chacune de
ces exceptions a sa justification, et aucune n'est absurde prise isolément. Leur accumulation l'est davantage. Près
d'un foyer sur trois bénéficie aujourd'hui d'une réduction ou d'un crédit d'impôt sur le revenu, pour 1\,520 euros en
moyenne\footnote{DGFiP, \emph{Statistiques} n\textsuperscript{o}~41, novembre 2025 : 12,8 millions de foyers bénéficiaires, 1\,522 euros en moyenne.}. Or chaque exception rétrécit l'assiette et
oblige à relever le taux sur le reste, ce qui alourdit le coût de l'ensemble. Celui qui ne coche aucune case paie pour
celui qui les coche toutes.

\textbf{Nous voulons que les décisions économiques dépendent de leur rentabilité réelle plutôt que de leur traitement
fiscal.} Une entreprise doit investir parce que son projet est rentable, et non parce que son comptable lui a trouvé un
dispositif ; un ménage doit choisir sa garde d'enfants, son placement ou son logement pour ce qu'ils lui apportent. Nos
propositions découlent de ce principe. Elles suppriment la quasi-totalité des niches pour appliquer des taux
raisonnables sur des assiettes larges, taxent moins le travail, l'investissement et la mobilité, et davantage la rente
foncière. Elles reconnaissent les charges de famille. %%C:p1 #1%% Nous ne construisons pas la fiscalité idéale des
manuels, qui taxerait la consommation plutôt que le revenu. Nous proposons une étape crédible dans cette direction.

\textbf{Le travail paie davantage.} Le nouvel impôt sur le revenu fusionne l'impôt sur le revenu, la CSG et la CRDS en un
seul impôt individuel, sans niches, dont le taux marginal monte sans à-coups jusqu'à 48~\%. Les pensions y sont imposées
comme les salaires et une franchise de 4\,700 euros par enfant remplace le quotient familial. La refonte est neutre pour
les finances publiques, la baisse des taux rendant aux contribuables ce que l'élargissement de l'assiette leur prend.

\textbf{L'impôt immobilier passe des transactions à la détention du sol.} Les droits de mutation sont divisés par deux.
Une taxe nationale sur la seule valeur du sol rapporte 15~Md€ et la taxe foncière communale est assise sur la valeur de
marché, en deux parts, à recette inchangée. Celui qui déménage paie moins, celui qui laisse inutilisé un terrain
constructible paie davantage.

\textbf{L'impôt sur les sociétés récompense l'investissement plutôt que la taille ou le secteur.} Le taux reste à 25~\%.
La déduction immédiate des investissements, portée par paliers à 100~\% en 2032, remplace le taux réduit des petites
entreprises et les crédits d'impôt sectoriels, et les déficits reportables sont revalorisés chaque année.

\textbf{La fiscalité du capital garde sa cohérence avec celle du travail.} Le prélèvement forfaitaire unique reste à
31,4~\%. Cumulé avec l'impôt sur les sociétés, il taxe un bénéfice distribué à 48,6~\%, presque exactement le taux
marginal maximal du nouvel impôt sur le revenu.

\section{2. Grands équilibres}

\textbf{Les quatre mesures dégagent 5,3~Md€ pour le solde public à terme, mais en coûtent 0,6 en 2032.} Les impôts
concernés rapportaient 381~Md€ en 2025. Ils en rapporteraient 387 à terme. La refonte de l'impôt sur le revenu est
neutre, la réforme foncière dégage un peu plus de 5~Md€ et l'impôt sur les sociétés, déficitaire pendant la montée en
charge de la déduction immédiate, s'équilibre presque une fois celle-ci achevée (tableau~\ref{tab:t1}).

\begin{table}[!htbp]\small\setlength{\tabcolsep}{4pt}
\caption{Recettes des impôts concernés et solde public (Md€ de 2025)}\label{tab:t1}
\begin{tabularx}{\linewidth}{@{}Xrrrr@{}}
\toprule
\textbf{Mesure et impôts concernés} & \textbf{Recettes 2025} & \textbf{Après réforme} & \textbf{Solde 2032} & \textbf{Solde à terme} \\
\midrule
\textbf{1. Nouvel impôt sur le revenu} & \textbf{260,2} & \textbf{260,2} & \textbf{0,0} & \textbf{0,0} \\
\quad Impôt sur le revenu, net des crédits d'impôt & 94,3 & fusionnés & & \\
\quad CSG & 156,6 & & & \\
\quad CRDS & 9,3 & & & \\
\textbf{2. Fiscalité foncière} & \textbf{61,8} & \textbf{67,3} & \textbf{+5,2} & \textbf{+5,4} \\
\quad Taxe foncière communale (sol et bâti) & 44,2 & 44,2 & 0,0 & 0,0 \\
\quad Taxe nationale sur la valeur du sol & 0,0 & 15,0 & +15,0 & +15,0 \\
\quad Droits de mutation à titre onéreux & 17,7 & 8,1 & $-$8,1 & $-$8,1 \\
\quad Déductibilité de la taxe sur le sol, calendrier & & & $-$1,8 & $-$1,5 \\
\textbf{3. Impôt sur les sociétés} & \textbf{59,3} & \textbf{59,2} & \textbf{$-$5,7} & \textbf{$-$0,1} \\
\textbf{4. Prélèvement forfaitaire unique (inchangé)} & & & \textbf{0,0} & \textbf{0,0} \\
\midrule \textbf{Total} & \textbf{381,3} & \textbf{386,6} & \textbf{$-$0,6} & \textbf{+5,3} \\
\bottomrule
\end{tabularx}
\par\vspace{2pt}{\raggedright\scriptsize\color{gris}Source : Insee, comptes nationaux 2025, tableau 3.217 (recettes 2025) ; calculs des auteurs. Après réforme et solde à terme : régime de croisière, une fois achevée la transition vers la déduction immédiate des investissements (recettes de 2040 rapportées à l'économie de 2025). Les droits de mutation de 2025 comprennent la hausse temporaire autorisée jusqu'en 2028 ; le solde est calculé par rapport au taux permanent (16,1 Md€). L'effet sur les prélèvements obligatoires est égal au solde public.\par}
\end{table}


\textbf{Les doublons avec la partie III sont retirés.} Le recentrage du crédit d'impôt recherche, la réduction d'impôt
pour mécénat, les zonages et le crédit d'impôt famille figurent parmi les aides que supprime la mesure 2 de la partie III.
La fin de l'effacement des plus-values à la transmission figure dans le paquet hauts patrimoines de sa mesure 3. Ces
recettes ne sont comptées qu'une fois, dans la partie III. Sans elles, la déduction immédiate des investissements n'est
financée qu'en partie.

\textbf{La transition pèse sur les années 2032 à 2035.} En passant à la déduction intégrale, l'État laisse les
entreprises déduire immédiatement leurs nouveaux investissements tout en continuant d'amortir les anciens. Il paie
ainsi deux fois pendant quelques années. Le solde des quatre mesures tombe à $-$7,7~Md€ en 2033, puis redevient positif
en 2036. Cumulé depuis 2028, il reste négatif de 2034 à 2036, au plus bas de 2,8~Md€, et atteint 17~Md€ en 2040.

\graphique{figure_bilan}{En 2032, l'impôt sur le revenu rend tout ce qu'il prélève et la marge du foncier finance la montée en charge de la déduction immédiate}{fig:bilan}

\BILANIR

\textbf{À terme, les prélèvements obligatoires augmentent de 0,2 point de PIB.} La marge de 5,3~Md€ s'ajoute à celle
de la partie III. Elle peut être rendue sous forme de baisses d'impôt, par exemple dans le barème du nouvel impôt sur le
revenu, ou financer une partie de la baisse des impôts de production.

\textbf{Trois arbitrages pèsent sur ces résultats.}
\begin{itemize}[leftmargin=12pt,itemsep=1pt,topsep=1pt]
\item \textbf{Qui finance la déduction immédiate\,?} La partie III divise par trois le coût du crédit d'impôt recherche et
en garde le gain. Supprimer aussi le tiers restant rapporterait 2,1~Md€ par an après réaction des entreprises et
porterait le solde de 2032 à +1,5~Md€. S'en tenir à une déduction de 50~\% diviserait par deux son coût permanent, au
prix d'une incitation plus faible.
\item \textbf{Faut-il baisser le prélèvement forfaitaire unique\,?} Le ramener à 30~\% coûterait 1,1~Md€, à 25~\% 3,8~Md€.
À 25~\%, un bénéfice distribué serait moins taxé qu'un salaire élevé (tableau~\ref{tab:t5}).
\item \textbf{Jusqu'où taxer le sol\,?} Chaque 0,05 point de taxe nationale rapporte environ 3~Md€. Au taux retenu, la
valeur des terrains baisse d'environ 7,5~\%, et la limite tient moins au rendement qu'au choc subi par les propriétaires
récents.
\end{itemize}

\section{3. Détail des chiffrages}

\subsection{3.1 Travailler davantage doit rapporter davantage}

\CHAPEAU{IR}

%%C:p1 #3%%

%%C:p1 #4%%

\graphique{g2_taux_marginaux}{Le taux marginal monte en ligne droite, sans marche, et ne dépasse jamais 48~\%}{fig:g2}

%%C:p2 #1%%

%%C:p2 #2%%

%%C:p2 #3%%

%%C:p2 #4%%

\graphique{g14_menages}{Une fusion « froide » favoriserait les célibataires ; la franchise par enfant fait des familles les premières gagnantes}{fig:g3}

%%C:p3 #3%%

%%C:p4 #1%%

\begin{table}[!htbp]\small\setlength{\tabcolsep}{4pt}
\caption{Composantes du nouvel impôt sur le revenu (Md€ de 2025, 2032 et à terme)}\label{tab:t2}
\begin{tabularx}{\linewidth}{@{}Xrc@{}}
\toprule
\textbf{Composante} & \textbf{Md€} & \textbf{Statut} \\
\midrule
Suppression des niches (réductions et crédits, exonérations, demi-parts) & +19,2 & DED \\
Pensions imposées comme les salaires (abattement, CSG, âge) & +15,6 & DED \\
Forfait de 500 € au lieu de l'abattement de 10 % et des frais réels & +25,7 & DED \\
Fin de la CSG déductible & +24,1 & DED \\
Impôt individuel, fin du quotient conjugal & +12,1 & DED \\
Fin du quotient familial & +16,2 & DED \\
Franchise de 4\,700 € par enfant & $-$15,1 & HYP \\
Baisse de tous les taux (nouveau barème) & $-$97,1 & CAL \\
Minimum vieillesse, contribution sur les hauts revenus & $-$0,8 & DED \\
\midrule \textbf{Solde public} & \textbf{0,0} & \\
\bottomrule
\end{tabularx}
\par\vspace{2pt}{\raggedright\scriptsize\color{gris}Source : modèle de foyers calé sur les statistiques de l'impôt sur le revenu 2024 (DGFiP) ; OpenFisca-France ; coûts des niches : Évaluation des voies et moyens (PLF 2026). Décomposition séquentielle avant réaction des comportements : l'ordre change la taille de chaque ligne, pas le total. Statut : DED, déduit d'une donnée observée ; CAL, calé (ici sur la neutralité par rapport au droit de 2032) ; HYP, paramètre choisi.\par}
\end{table}


\textbf{Le chiffrage est statique et calé sur le droit de 2032.} Les composantes du tableau~\ref{tab:t2} viennent d'un
modèle de foyers calé sur les déclarations de revenus de 2024. Les taux du barème sont calés pour que la recette soit
celle du droit en vigueur en 2032, une fois éteints les dispositifs qui doivent disparaître d'ici là. La franchise par
enfant est le principal paramètre de choix politique. Chaque millier d'euros de franchise coûte de l'ordre de 3~Md€, à
compenser par les taux. Deux effets favorables ne sont pas retenus, le surcroît d'activité que permet la baisse des taux
marginaux, de 1 à 5~Md€ selon les élasticités, et l'interaction avec la mesure 1 de la partie III. Celle-ci relève le
salaire net en supprimant des cotisations, et le nouvel impôt, assis sur le revenu après cotisations, en reprendrait une
part plus grande que l'impôt actuel.

\subsection{3.2 Une fiscalité foncière au service de la croissance}

\CHAPEAU{FONCIER}

%%C:p4 #3%%

\graphiqueetroit{g15_foncier}{À recette nette plus élevée, l'impôt immobilier passe des transactions à la détention du sol}{fig:g15}{0.44}

%%C:p4 #4%%

\deuxgraphiques{g10_foncier_pib}{La valeur du sol bâti a plus que quadruplé rapportée au PIB depuis la fin des années 1990}{fig:g10}%
{g17_construire}{La taxe sur le sol renchérit le terrain laissé vide sans alourdir l'impôt dû pour avoir construit}{fig:g6}

%%C:p5 #1%%

\begin{table}[!htbp]\small\setlength{\tabcolsep}{4pt}
\caption{Composantes de la réforme foncière (Md€ de 2025)}\label{tab:t3}
\begin{tabularx}{\linewidth}{@{}Xrrc@{}}
\toprule
\textbf{Composante} & \textbf{2032} & \textbf{À terme} & \textbf{Statut} \\
\midrule
Taxe nationale sur la valeur du sol (taux de 0,24 %) & +15,0 & +15,0 & HYP \\
Déductibilité de la taxe sur le sol, calendrier de recouvrement & $-$1,8 & $-$1,5 & DED \\
Droits de mutation divisés par deux & $-$8,1 & $-$8,1 & OBS \\
Taxe communale sur la valeur de marché, en deux parts & 0,0 & 0,0 & HYP \\
\midrule \textbf{Solde public} & \textbf{+5,2} & \textbf{+5,4} & \\
\bottomrule
\end{tabularx}
\par\vspace{2pt}{\raggedright\scriptsize\color{gris}Source : valeurs foncières des ventes (DVF), comptes de patrimoine de l'Insee ; droits de mutation au taux permanent, 16,1 Md€ en 2025 (OFGL, DGFiP) ; calculs des auteurs. Assiette du sol : 6\,156 Md€ après la baisse de valeur que provoque la taxe (construite, non observée). La taxe communale garde sa recette de 2025.\par}
\end{table}


\textbf{L'assiette du sol est construite et non observée.} La valeur du sol vient des prix des ventes et des comptes de
patrimoine de l'Insee. Elle intègre la baisse de valeur que la taxe provoque, d'environ 7,5~\%. Chaque point de droits de
mutation représente 2,8~Md€. Le surcroît de ventes qu'entraînerait la baisse des droits n'est pas compté.

\subsection{3.3 Une fiscalité des entreprises au service de l'investissement}

\CHAPEAU{IS}

%%C:p6 #2%%

\begin{table}[!htbp]\small\setlength{\tabcolsep}{4pt}
\caption{Composantes de la réforme de l'impôt sur les sociétés (Md€ de 2025)}\label{tab:t4}
\begin{tabularx}{\linewidth}{@{}Xrrc@{}}
\toprule
\textbf{Composante} & \textbf{2032} & \textbf{À terme} & \textbf{Statut} \\
\midrule
Fin du taux réduit de 15 % des petites entreprises & +2,0 & +2,0 & DED \\
Fin des crédits d'impôt culture (cinéma, jeux vidéo, spectacle) & +0,4 & +0,5 & DED \\
Déduction immédiate des investissements (25 % en 2029, 50 % en 2030, 100 % en 2032) & $-$7,0 & $-$1,5 & DED \\
Revalorisation des déficits reportables au taux de l'État & $-$1,1 & $-$1,1 & HYP \\
\midrule \textbf{Solde public} & \textbf{$-$5,7} & \textbf{$-$0,1} & \\
\midrule \emph{Pour mémoire, compté en partie III : recentrage du crédit d'impôt recherche} & \emph{+5,2} & \emph{+5,2} & \\
\emph{Pour mémoire, compté en partie III : mécénat, zonages, crédit famille} & & & \\
\emph{Variante : suppression du tiers restant du crédit d'impôt recherche} & \emph{+2,1} & \emph{+2,1} & HYP \\
\bottomrule
\end{tabularx}
\par\vspace{2pt}{\raggedright\scriptsize\color{gris}Source : Insee (recettes 2025) ; DGFiP, \emph{Statistiques} n\textsuperscript{o}~51 ; Évaluation des voies et moyens (PLF 2026) ; calculs des auteurs. Rendements après une décote de 10 à 25~\% pour la réaction des entreprises. La déduction immédiate coûte davantage en 2032 qu'à terme parce que l'État rembourse encore les amortissements des investissements passés. Variante : 2,6 Md€ après recentrage, décote de 20~\%.\par}
\end{table}


\textbf{Deux paramètres fragiles déterminent ce bloc.} Le coût de la revalorisation des déficits dépend du volume de
déficits imputés chaque année, qui n'est pas publié. Il varie de 0,4 à 2,2~Md€. Le coût de la déduction immédiate dépend
du rythme de l'investissement, supposé suivre le PIB. Un effet favorable n'est pas retenu, comme dans la partie III. La
baisse des impôts de production élargit le bénéfice imposable et relève d'autant le rendement de l'impôt sur les sociétés.

\subsection{3.4 Préserver les acquis du prélèvement forfaitaire unique}

\CHAPEAU{PFU}

%%C:p7 #2%%

\PFUFIN

\begin{table}[!htbp]\small\setlength{\tabcolsep}{4pt}
\caption{Coût d'une baisse du prélèvement forfaitaire unique (Md€ de 2025, option non retenue)}\label{tab:t5}
\begin{tabularx}{\linewidth}{@{}Xrrrr@{}}
\toprule
\textbf{Option} & \textbf{Statique} & \textbf{Après réaction} & \textbf{Fourchette} & \textbf{Taux cumulé} \\
\midrule
Prélèvement forfaitaire unique à 30~\% & $-$1,5 & $-$1,1 & $-$1,4 à $-$1,0 & 47,5~\% \\
Prélèvement forfaitaire unique à 25~\% & $-$5,2 & $-$3,8 & $-$4,8 à $-$3,4 & 44~\% \\
\bottomrule
\end{tabularx}
\par\vspace{2pt}{\raggedright\scriptsize\color{gris}Source : DGFiP (dividendes, intérêts et plus-values déclarés au titre de 2024, 80 Md€) ; élasticité du revenu déclaré de 0,22 à 0,96 (Lefebvre, Lehmann et Sicsic) ; Sénat pour le rendement constaté de la hausse de 2026. Taux cumulé : impôt sur les sociétés à 25~\% puis prélèvement sur le dividende, à comparer aux 48~\% du haut du nouvel impôt sur le revenu.\par}
\end{table}


\subsection{3.5 Statut des données et suites}

\begin{itemize}[leftmargin=12pt,itemsep=1pt,topsep=1pt]
\item \textbf{Données observées.} Recettes 2025 par impôt (Insee, tableau 3.217), déclarations de revenus de 2024 et
frais réels déclarés (DGFiP), coûts des niches (Évaluation des voies et moyens, PLF 2026), droits de mutation 2025 (OFGL,
DGFiP), bénéfices imposés au taux normal (DGFiP, \emph{Statistiques} n\textsuperscript{o}~51), revenus du capital déclarés.
\item \textbf{Données construites.} Valeur du sol et sa baisse sous l'effet de la taxe, répartition des foyers et des
frais professionnels, profil de l'investissement et des amortissements, volume des déficits reportables imputés,
réaction des entreprises à la fin des niches.
\item \textbf{Suites proposées.} Obtenir un échantillon de déclarations de la DGFiP pour un bilan par décile de niveau de
vie, chiffrer l'interaction du nouvel impôt avec la mesure 1 de la partie III, actualiser la valeur du sol, arbitrer le
financement de la déduction immédiate.
\end{itemize}

\section{4. Pour aller plus loin : cesser de pénaliser la construction et les revenus irréguliers}

\textbf{Deux chantiers, hors du scénario chiffré, prolongent la même logique : que l'impôt cesse de décourager ceux
qui construisent et ceux dont les revenus sont irréguliers.}

\textbf{Notre réforme doit inciter non seulement les propriétaires à construire, mais aussi les maires à les laisser
construire.} Depuis la suppression de la taxe d'habitation, un logement neuf ne rapporte plus à sa commune que la taxe
foncière, et seulement après deux ans d'exonération. Le maire supporte en revanche le coût des écoles, des réseaux et,
souvent, l'hostilité des riverains. Notre réforme aggrave ce déséquilibre : parce qu'elle taxe le sol plus que le bâti,
la recette annuelle qu'un appartement neuf de 290\,000 euros apporte à sa commune tomberait d'environ 780 à 310 euros.
Une dotation automatique de 5\,000 euros par logement supplémentaire, net des démolitions, versée dans les zones
tendues sans dossier ni appel à projets, coûterait 0,75 à 1,5 milliard d'euros par an. Elle serait prise sur la dotation
globale de fonctionnement à enveloppe constante, de sorte que les communes qui ne construisent pas y perdraient.
L'Angleterre a pratiqué une prime de ce type, dont l'évaluation n'a pas démontré d'effet sur la construction. Nous n'en
chiffrons donc aucun.

\textbf{À revenu cumulé identique, notre fiscalité pénalise celui qui prend le risque d'avoir des revenus irréguliers.}
Avec notre barème, Antoine, qui gagne 50\,000 euros par an pendant cinq ans, paie 47\,100 euros d'impôt. Béatrice, qui
lance son activité et gagne successivement 10\,000, 20\,000, 30\,000, 70\,000 puis 120\,000 euros, soit le même total,
paie 58\,400 euros\footnote{Calculs des auteurs avec le barème proposé, impôt calculé chaque année sur le revenu de l'année.}. La progressivité, juste lorsqu'elle compare des revenus durablement différents, pénalise ici une
même richesse répartie inégalement dans le temps. Le droit actuel ne corrige ce défaut que pour quelques professions,
agriculteurs, artistes, écrivains ou sportifs. Nous proposons de remplacer ces régimes par une option ouverte à tous les
revenus professionnels : pendant cinq ans, chacun paierait l'impôt dû sur son revenu moyen, régularisé chaque année et à
la sortie du dispositif, en cas de cessation d'activité, de décès ou de départ à l'étranger. L'écart entre Antoine et
Béatrice disparaîtrait. Son coût ne peut être chiffré sans suivre les mêmes contribuables pendant plusieurs années, mais
les régimes actuels coûtent peu.
